\section*{Tóm tắt}
\addcontentsline{toc}{section}{Tóm tắt}

X-quang ngực (CXR) là xét nghiệm hình ảnh phổ biến nhất, và nhiều ca ung thư phổi lần đầu lộ ra dưới dạng
nốt hoặc khối trên CXR, nhưng nốt nhỏ và mờ rất dễ bị bỏ sót. Các mô hình ngôn ngữ--thị giác (VLM) đang được
giới thiệu như trợ lý đọc CXR, song các đánh giá độc lập năm 2026 cho thấy sensitivity zero-shot với nốt phổi
chỉ đạt 0,6--27,3\%; MedGemma-4B đạt sensitivity 5,2\% ở specificity 100\% --- một \emph{chế độ báo bình
thường}. Trong khi đó, hệ thống CAD dựa trên CNN đạt AUROC 0,92--0,99. Đồ án xây dựng và đánh giá một
\textbf{nguyên mẫu trợ lý AI} cho biết ảnh CXR có nốt/khối phổi hay không và chỉ ra vị trí bằng bounding box,
dựa trên VLM cỡ nhỏ tinh chỉnh LoRA.

Đồ án gồm ba phần. (i) \textbf{Benchmark chuẩn hoá}: prompt và bộ phân tích đầu ra công khai, đầu ra JSON có
cấu trúc, ngưỡng chọn trên tập validation và chạy lặp; so sánh MedGemma-4B và Qwen2.5-VL-3B với hai baseline
DenseNet-121 và YOLOv8 trên VinBigData (trong phân phối), NIH ChestX-ray14 và JSRT (ngoài phân phối, có phân
cấp độ mờ). (ii) \textbf{Tinh chỉnh LoRA} MedGemma-4B cho bài toán chỉ phân loại và phân loại kèm định vị, có
và không có cân bằng lớp, nhằm kéo mô hình ra khỏi chế độ báo bình thường. (iii) \textbf{Kiểm tra trợ lý có
nhìn đúng chỗ không}: IoU định vị, so sánh che nốt với che vùng ngẫu nhiên cùng diện tích, và phân tầng theo
kích thước và độ mờ của nốt. Chỉ số chính là sensitivity ở specificity 0,90, với mục tiêu 0,60; các chỉ số phụ
gồm AUC-ROC, AUC-PR, calibration và FROC, kèm khoảng tin cậy 95\% bằng bootstrap.

\textbf{Kết quả.} \textcolor{badred}{[Điền sau thí nghiệm E1--E5.]}

\textbf{Kết luận.} \textcolor{badred}{[Điền sau thí nghiệm E1--E5.]} Đây là nguyên mẫu nghiên cứu, không dùng
cho quyết định lâm sàng.

\vspace{0.6em}
\noindent\textbf{Từ khoá:} trợ lý AI, mô hình ngôn ngữ--thị giác, LoRA, nốt phổi, X-quang ngực, phát hiện và
định vị tổn thương, mất cân bằng lớp.
